\section{Methodology}
\label{sec:methodology}

\paragraph{Reranking setup.}
Given a query $q$ and a first-stage BM25 run
~\citep{robertson2009probabilistic}, let
$P=(d_1,\ldots,d_N)$ denote the top-$N$ candidate pool in BM25 order,
with associated BM25 scores. The task is to produce a full reranking
$R=(r_1,\ldots,r_N)$ of this pool; evaluation metrics such as nDCG@10 are
then computed from the prefix of $R$. This paper studies the matched
operating point $k=N$ for all methods and, for MaxContext, an initial
prompt window $W=N$: the whole reranking pool is visible in the first
LLM call. We use passage length $\ell=512$ for rendering candidate
passages. After rendering, a runtime tokenization preflight checks that
the full prompt fits the model context window; prompts are not truncated.

\paragraph{Compared methods.}
Our method axis contains seven rerankers. Four are standard setwise
baselines following the local-window setwise formulation of
\citet{zhuang2024setwise}: TopDown setwise with bubblesort, TopDown
setwise with heapsort, BottomUp setwise with bubblesort, and BottomUp
setwise with heapsort. These baselines use the public
\texttt{llm-rankers} implementation~\citep{ielab_llmrankers} and retain
the usual sorting-based scheduling of local setwise comparisons. The
remaining three methods are MaxContext variants: best-only
\textsc{MaxContext-TopDown}, worst-only \textsc{MaxContext-BottomUp}, and
joint best--worst \textsc{MaxContext-DualEnd}. Standard DualEnd sorting
baselines are not part of the EMNLP method matrix.

\paragraph{MaxContext skeleton.}
MaxContext maintains a live pool $L$, initialized to the BM25 top-$N$
list $P$. In each LLM call, the prompt contains the query and every
document currently in $L$. Documents are labeled with local numeric
identifiers $1,\ldots,|L|$, which replace the legacy single-letter
interface and allow the full live pool to be addressed when
$|L|$ exceeds the letter-label range. The relative order of unselected
documents is preserved after each removal, so the only changes to $L$
come from documents that have already been placed into the output
ranking.

The model output is strictly parsed as integer labels. For a
single-extreme call, the parser accepts exactly one integer in
$[1,|L|]$. For a dual-end call, it accepts exactly two distinct integers
in $[1,|L|]$, one for the best document and one for the worst document.
Out-of-range labels, duplicate labels, missing labels, extra labels, or
refusal-like outputs are treated as parse failures. We do not repair such
outputs with fallback heuristics; the run is invalidated and recorded by
the parser telemetry. Thus, every successful MaxContext run returns a
true permutation of the original BM25 pool.

\paragraph{Three MaxContext variants.}
\textsc{MaxContext-TopDown} asks the model to select the best document in
the live pool. The selected document is written to the next free head
position of $R$ and removed from $L$. \textsc{MaxContext-BottomUp} asks
the model to select the worst document in the live pool. The selected
document is written to the next free tail position of $R$ and removed
from $L$. Both single-extreme variants stop when two documents remain;
the final two-document residual is ordered deterministically by the
original BM25 run, with the higher-BM25 document placed above the
lower-BM25 document.

\textsc{MaxContext-DualEnd} asks for both extremes in the same call: the
best document and the worst document in the current live pool. The best
document is placed at the next free head position, the worst document is
placed at the next free tail position, and both are removed from $L$.
Algorithm~\ref{alg:dualend} formalizes this variant. The two
single-extreme variants are obtained by dropping one of the two
placements and using the two-document BM25 endgame described above.

\begin{algorithm}[t]
\small
\caption{\textsc{MaxContext-DualEnd}}
\label{alg:dualend}
\begin{algorithmic}[1]
\State \textbf{Input:} query $q$, ordered BM25 top-$N$ pool $P=(d_1,\ldots,d_N)$
\State \textbf{Output:} full reranking $R$ of length $N$
\State $L \gets P$
\State $R \gets$ empty array of length $N$
\State $\texttt{top} \gets 1$; \ $\texttt{bot} \gets N$
\While{$|L| \ge 2$}
  \State $y \gets \textsc{LLM}_{\mathrm{dual}}(q,L)$
  \State parse $y$ as $(i_{\mathrm{best}}, i_{\mathrm{worst}})$
  \Comment{distinct labels in $[1,|L|]$}
  \If{the parse is invalid}
    \State \textbf{abort}
  \EndIf
  \State $R[\texttt{top}] \gets L[i_{\mathrm{best}}]$
  \State $R[\texttt{bot}] \gets L[i_{\mathrm{worst}}]$
  \State $\texttt{top} \gets \texttt{top}+1$
  \State $\texttt{bot} \gets \texttt{bot}-1$
  \State delete $L[\max(i_{\mathrm{best}}, i_{\mathrm{worst}})]$
  \State delete $L[\min(i_{\mathrm{best}}, i_{\mathrm{worst}})]$
\EndWhile
\If{$|L| = 1$}
  \State $R[\texttt{top}] \gets L[1]$
\EndIf
\State \textbf{return} $R$
\end{algorithmic}
\end{algorithm}

\paragraph{LLM-call counts.}
At the matched MaxContext operating point with initial $W=N$ and output
depth $k=N$, \textsc{MaxContext-DualEnd} removes two documents per LLM
call and therefore uses
\[
C_{\mathrm{DualEnd}}(N)=\left\lfloor \frac{N}{2} \right\rfloor
\]
LLM calls. Each single-extreme MaxContext variant removes one document
per LLM call until two documents remain, so for the pool sizes studied in
this paper,
\[
C_{\mathrm{TopDown}}(N)=C_{\mathrm{BottomUp}}(N)=N-2,
\]
plus one deterministic two-document BM25 endgame that does not call the
LLM. At $N=100$, this gives $50$ LLM calls for
\textsc{MaxContext-DualEnd} versus $98$ LLM calls for either
single-extreme MaxContext control. The ratio is $50/98 \approx 0.5102$
and approaches $1/2$ from above as $N$ grows.

This call-count comparison is deterministic but deliberately narrow. It
is a claim about serial LLM invocations within the MaxContext family, not
a claim that DualEnd always uses half the tokens, half the FLOPs, or half
the wall-clock time. Prompt lengths differ across variants because the
live pool shrinks at different rates, and standard setwise baselines have
sorting-dependent call schedules. We therefore report LLM calls,
wall-clock time, input tokens, output tokens, and total tokens separately
in the experiments.

\paragraph{Hard invariants.}
The EMNLP configuration fixes the following invariants for MaxContext:
initial $W=N$, output depth $k=N$, pool sizes $N \le 100$, passage length
$\ell=512$, numeric one-based labels, strict context-fit preflight,
strict no-truncation after rendering, strict no-parse-fallback, and no
permutation self-consistency. MaxContext is run only with generation-based
causal LLMs from the Qwen3.5, Llama-3.1, and Ministral-3 families.
Likelihood-scoring T5-style rerankers are excluded because the
joint best--worst objective does not reduce to an independent
single-label likelihood score without changing the method being tested.